\subsection{紧空间上连续函数代数的子代数}

设 X 为紧拓扑空间，B 为 \(\mathcal{C}(\mathrm{X})\) 的一个幺子代数。我们用 ev 表示映射 \(x \mapsto ev_{x}\) ，它从 X 到 \(\mathrm{X}(\mathrm{B})\) 中，且满足 \(\mathrm{ev}_{x}(f) = f(x)\) 对所有 \(f \in B\) 成立。我们用 j 表示 B 到 \(\mathcal{C}(\mathrm{X})\) 中的单射，于是 \(\mathrm{X}(j) = \mathrm{ev}\) 。

命题 1.　设 \(f \mapsto \|f\|_{\mathrm{B}}\) 是一个范数，它使 B 具有 Banach 代数的结构。

\begin{enumerate}
\def\labelenumi{\alph{enumi})}
\item
  映射 j 在 \(\mathcal{L}(B; \mathcal{C}(\mathrm{X}))\) 中的范数 ≤ 1；
\item
  代数 B 的根为零；
\item
  映射 ev 从 X 到 X(B) 中是连续的；
\item
  若 B 分离 X 的点，则映射 ev 是从 X 到 X(B) 的一个闭子集上的同胚。
\end{enumerate}

可以把 X 等同于 \(\mathsf{X}(\mathcal{C}(\mathsf{X}))\) ，并把 \(\mathcal{G}_{\mathcal{C}(\mathsf{X})}\) 等同于恒等映射（I，第 36 页，例 1）。于是映射 ev 等同于 \(\mathsf{X}(j)\) ，这就证明了 c)。

对每个函数 \(f \in \mathrm{B}\) 与每个 \(x \in \mathrm{X}\) ，有 \(\mathcal{G}_{\mathrm{B}}(f)(\mathrm{ev}_x) = f(x)\) ，由此得

\[
\| f \| _ {\mathcal {C} (\mathrm{X})} = \sup _ {x \in \mathrm{X}} | \mathcal {G} _ {\mathrm{B}} (f) (\mathrm{ev} _ {x}) | \leqslant \varrho_ {\mathrm{B}} (f) \leqslant \| f \| _ {\mathrm{B}},
\]

（参见 I，第 38 页，命题 7，a)）由此便得 \(a\) ）。此外，这表明 B 的 Gelfand 变换是单射，从而得 \(b\) ）（I，第 38 页，命题 8）。若 B 分离 X 的点，则映射 ev 是单射，由于 X 是紧的，由此得 \(d\) ）。

现在转向映射 ev 的满射性。

命题 2.　设 \(f \mapsto \|f\|_{\mathrm{B}}\) 是一个范数，使得 \((\mathrm{B}, \| \cdot \|_{\mathrm{B}})\) 是一个 Banach 代数。

\begin{enumerate}
\def\labelenumi{\alph{enumi})}
\item
  若映射 ev 是满射，则代数 B 是 \(\mathcal{C}(\mathrm{X})\) 的全子代数；
\item
  假设 B 是 \(\mathcal{C}(\mathrm{X})\) 的全子代数，且存在元素 \(a \in \mathrm{B}\) ，使得诸元素 \(f(a)\) （其中 \(f\) 取遍 \(\mathbf{C}\) 上在 \(\mathrm{Sp}_{\mathrm{B}}(a)\) 上无极点的有理函数）所成的集合在 B 中稠密。则映射 ev 是满射；
\item
  若 B 是 \(\mathcal{C}(\mathrm{X})\) 的全对合子代数，则映射 ev 是满射。
\end{enumerate}

断言 \(a\) ）由 I，第 40 页，命题 10 得出，断言 \(b\) ）由 I，第 41 页，命题 11 得出，因为 B 的由 \(a\) 生成的全子代数就是诸元素 \(f(a)\) （其中 \(f\) 取遍 \(\mathbf{C}\) 上在 \(\mathrm{Sp}_{\mathrm{B}}(a)\) 上无极点的有理函数）所成的集合（I，第 6 页，引理及 I，第 37 页，命题 6，\(b\) ））。

我们证明 \(c\) ）。于是假设 B 是 \(\mathcal{C}(\mathrm{X})\) 的全对合子代数。为证明 ev 是满射，只需证明：对每个 \(\chi \in \mathrm{X}(\mathrm{B})\) ，存在 \(y \in \mathrm{X}\) ，使得 \(\operatorname{Ker}(\chi) = \operatorname{Ker}(\mathrm{ev}_y)\) （I，第 30 页，定理 2）。设 \(\mathrm{I} = \operatorname{Ker}(\chi)\) 。这是 B 的一个极大理想。设 \(\Phi\) 为这样的 \(x \in \mathrm{X}\) 全体所成的集合：\(f(x) = 0\) 对所有 \(f \in \mathrm{I}\) 成立。我们证明 \(\Phi\) 非空。若不然，由于 X 是紧的，就会存在整数 \(n \geqslant 1\) 、X 的开覆盖 \((\mathrm{V}_1, \ldots, \mathrm{V}_n)\) ，并且对每个整数 \(i\) （\(1 \leqslant i \leqslant n\)），存在函数 \(f_i \in \mathrm{I}\) ，使得 \(f_i(x) \neq 0\) 对所有 \(x \in \mathrm{V}_i\) 成立。由于代数 B 是 \(\mathcal{C}(\mathrm{X})\) 的对合子代数，函数

\[
f = \sum_ {i = 1} ^ {n} f _ {i} \overline {{f}} _ {i}
\]

将属于 I。但 \(f(x) > 0\) 对所有 \(x \in X\) 成立，因此 f 在 \(\mathcal{C}(X)\) 中可逆。由于假设 B 是 \(\mathcal{C}(X)\) 的全子代数，函数 \(f \in I\) 就会在 B 中可逆，而这是不可能的。因此，集合 \(\Phi\) 非空。设 y 是 \(\Phi\) 的一个元素；特征标 \(ev_{y}\) 的核包含 I，从而等于 I。

例.　设 X = \([0,1]\) ，\(n \geqslant 0\) 为整数。设 B 是由这样的函数 \(f: X \to C\) 组成的代数：它们在 \([0,1]\) 上具有直到 n 阶的连续导数，且赋以 I，第 18 页，例 4 中所考虑的范数。则 B 是 \(\mathcal{C}(X)\) 的分离 X 的点的全对合子代数，因而 \(\mathsf{X}(\mathsf{B})\) 可等同于 X。

我们考虑定义在 \(R_{+}\) 上、取值于 \(R \cup \{-\infty\}\) 中的对数函数，其定义为 \(\log(0) = -\infty\) 。

命题 3.　设 X 为紧空间。设 B 为 \(\mathcal{C}(\mathrm{X})\) 的幺 Banach 子代数，赋以诱导范数，且分离 X 的点。我们把 X 等同于 \(\mathrm{X(B)}\) 的一个闭子集（命题 1，d)）。

\begin{enumerate}
\def\labelenumi{\alph{enumi})}
\item
  对每个 \(f \in \mathbf{B}\) ，Gelfand 变换 \(\mathcal{G}_{\mathrm{B}}(f)\) 是 \(\mathsf{X}(\mathbf{B})\) 上的连续函数，它是 \(f\) 的扩张，且满足 \(\| f \| = \sup |\mathcal{G}_{\mathrm{B}}(f)|\) 。特别地，\(\mathcal{G}_{\mathrm{B}}\) 是从 B 到 \(\mathcal{C}(\mathsf{X}(\mathbf{B}))\) 的一个 Banach 子代数上的等距同构；
\item
  设 B* 是 B 的可逆元素所成的集合。对每个特征标 \(\chi \in \mathsf{X}(\mathsf{B})\) ，存在 X 上的质量为 1 的正测度 \(\mu\) ，使得对每个 \(f \in \mathsf{B}^*\) ，有
\end{enumerate}

\[
\log (| \chi (f) |) = \int_ {\mathrm{X}} \log (| f |) d \mu .
\]

此外，对每个 \(f \in \mathrm{B}\) ，有

\[
\chi (f) = \int_ {\mathrm{X}} f d \mu ;
\]

\begin{enumerate}
\def\labelenumi{\alph{enumi})}
\setcounter{enumi}{2}
\tightlist
\item
  我们假设 \(\mathcal{C}_{\mathbf{R}}(\mathrm{X})\) 的每个元素都是 B 中函数的实部的一致极限。则对每个 \(\chi \in \mathrm{X}(\mathrm{B})\) ，存在 X 上唯一的测度 \(\mu_{\chi} \geqslant 0\) ，使得对每个函数 \(f \in B\) ，有
\end{enumerate}

\[
\chi (f) = \int_ {\mathrm{X}} f d \mu_ {\chi}.
\]

此外，对每个函数 \(f \in \mathrm{B}\) ，有

\[
\log (| \chi (f) |) \leqslant \int_ {\mathrm{X}} \log (| f |) d \mu_ {\chi};
\]

由于函数 \(\log|f|\) 上有界，该积分在 \(R \cup \{-\infty\}\) 中存在。

断言 \(a\) ）由把 X 等同于 \(\mathsf{X}(\mathsf{B})\) 的闭子集以及下列不等式得出：

\[
\| f \| = \sup _ {x \in \mathrm{X}} | f (x) | \leqslant \sup _ {\chi \in \mathrm{X} (\mathrm{B})} | \chi (f) | = \sup | \mathcal {G} _ {\mathrm{B}} (f) | = \varrho_ {\mathrm{B}} (f) \leqslant \| f \|
\]

对所有 \(f \in B\) 成立。

设 \(\chi \in \mathsf{X}(\mathbf{B})\) ，\(n\) 为正整数。设 \(\lambda_1, \ldots, \lambda_n \in \mathbf{R}\) 与 \(f_1, \ldots, f_n \in \mathbf{B}^*\) 。这时我们有

\[
\sum_ {i = 1} ^ {n} \lambda_ {i} \log (| \chi (f _ {i}) |) \leqslant \sup _ {x \in X} \left(\sum_ {i = 1} ^ {n} \lambda_ {i} \log (| f _ {i} (x) |)\right).\tag{1}
\]

事实上，由连续性，只需在实数 \(\lambda_{i}\) 为有理数时证明这个不等式。通分后，可归结为对所有 i 有 \(\lambda_{i} \in Z\) 的情形。这时不等式变为

\[
\log \left(\left| \chi \left(f _ {1} ^ {\lambda_ {1}} \dots f _ {n} ^ {\lambda_ {n}}\right) \right|\right) \leqslant \sup _ {x \in X} \log \left(\left| \left(f _ {1} ^ {\lambda_ {1}} \dots f _ {n} ^ {\lambda_ {n}}\right) (x) \right|\right),
\]

并由 \(\|\chi\|=1\) 这一事实（I，第 29 页，定理 1）得出。

设 B′ 为 \(\mathcal{C}_{\mathbf{R}}(\mathrm{X})\) 的由诸函数 \(\log(|f|)\) （\(f \in B^{*}\)）生成的向量子空间。估计式 (1) 证明：存在 B′ 上范数 \(\leqslant 1\) 的线性型 h，使得 \(\log(|\chi(f)|) = h(\log(|f|))\) 对所有 \(f \in B^{*}\) 成立。由 Hahn–Banach 定理（EVT，II，第 24 页，推论 2），线性型 h 可扩张为线性型 \(\mu\) ，它有范数 \(\leqslant 1\) ，定义在 \(\mathcal{C}_{\mathbf{R}}(\mathrm{X})\) 上，即 X 上的实测度 \(\mu\) ，它满足 \(\|\mu\| \leqslant 1\) （INT，III，§1，第 5 节）。取常数 \(e = \exp(1)\) 作为 B* 的元素 f，便看出 \(1 = \mu(1)\) 。于是，把 \(\mu = \mu^{+} - \mu^{-}\) 写成两个互奇异正测度之差（INT，III，§1，第 6 节，定理 3），便得

\[
1 = \mu^ {+} (1) - \mu^ {-} (1) \leqslant \mu^ {+} (1) + \mu^ {-} (1) = \| \mu \| \leqslant 1,
\]

由此得 \(\mu = \mu^{+} \geqslant 0\) 且 \(\|\mu\| = 1\) 。

对每个 \(f \in \mathrm{B}\) ，有 \(\exp(f) \in \mathrm{B}^*\) ，从而

\[
\begin{array}{r l} \int_ {\mathrm{X}} \mathcal {R} (f) d \mu = \int_ {\mathrm{X}} \log (| \exp (f) |) d \mu = \log (| \chi (\exp (f)) |) \\ & = \log (| \exp (\chi (f)) |) = \mathcal {R} (\chi (f)), \end{array}
\]

这里我们用到了 I，第 66 页，推论 1。把这个等式应用于 \(if\) ，便得出 \(\int_{\mathrm{X}} f d\mu = \chi(f)\) 。这就建立了 \(b\) ）。

我们置身于 \(c\) ）的假设之下。\(\mu_{\chi}\) 的存在性由 \(b\) ）得出。另一方面，我们有 \(\mu_{\chi}(\mathcal{R}(f)) = \mathcal{R}(\chi(f))\) 对所有 \(f \in B\) 成立。由于按假设，函数 \(f \in B\) 的实部在 \(\mathcal{C}_{\mathbf{R}}(\mathrm{X})\) 中稠密，测度 \(\mu_{\chi}\) 由 \(\chi\) 唯一确定。

设 \(f \in \mathrm{B}\) 。设 \(\varepsilon > 0\) 为实数。由假设，存在函数 \(g \in \mathrm{B}\) ，使得

\[
\mathcal {R} (g) - \varepsilon \leqslant \log (| f | + \varepsilon) \leqslant \mathcal {R} (g) + \varepsilon .\tag{2}
\]

设 \(h = \exp(g) \in \mathrm{B}^*\) 。由 (2)，我们有

\[
| h | e ^ {- \varepsilon} \leqslant | f | + \varepsilon \leqslant | h | e ^ {\varepsilon}.\tag{3}
\]

这个上界蕴涵 \(|fh^{-1}| \leqslant e^{\varepsilon}\) ，由此得 \(|\chi (fh^{-1})| \leqslant e^{\varepsilon}\) ，从而

\[
\log (| \chi (f) |) \leqslant \log (| \chi (h) |) + \varepsilon = \int_ {\mathrm{X}} \log (| h |) d \mu_ {\chi} + \varepsilon .\tag{4}
\]

于是 (3) 中的下界蕴涵

\[
\log (| \chi (f) |) \leqslant \int_ {\mathrm{X}} \log (| f | + \varepsilon) d \mu_ {\chi} + 2 \varepsilon .
\]

让 \(\varepsilon\) 趋于 0，我们推出

\[
\log (| \chi (f) |) \leqslant \int_ {\mathrm{X}} \log (| f |) d \mu_ {\chi}.
\]

证明完毕。

